
\documentclass[10pt,twocolumn,a4paper]{article}

\usepackage[T1]{fontenc}
\usepackage[utf8]{inputenc}
\usepackage{lmodern}
\usepackage[a4paper,margin=1.6cm,columnsep=0.6cm]{geometry}
\usepackage{microtype}
\usepackage{xcolor}
\usepackage{graphicx}
\usepackage{booktabs}
\usepackage{enumitem}
\usepackage{titlesec}
\usepackage[numbers,sort&compress]{natbib}
\usepackage[hidelinks]{hyperref}

% ---- compact layout -------------------------------------------------
\titleformat{\section}{\normalsize\bfseries}{\thesection}{0.5em}{}
\titlespacing*{\section}{0pt}{0.8ex plus 0.3ex}{0.4ex}
\setlist{nosep,leftmargin=*}
\setlength{\parskip}{0pt}
\pagestyle{empty}

% ---- placeholder macro (gray italic so leftovers are easy to spot) ---
\newcommand{\ph}[1]{\textcolor{gray}{\textit{[#1]}}}

% ---- compact bibliography -------------------------------------------
\renewcommand{\bibsection}{\section*{References}}
\setlength{\bibsep}{0pt}
\renewcommand{\bibfont}{\footnotesize}
\newcommand{\doi}[1]{\href{https://doi.org/#1}{doi:\,#1}}

% ---- title block ----------------------------------------------------
\makeatletter
\newcommand{\@affiliation}{}
\newcommand{\affiliation}[1]{\renewcommand{\@affiliation}{#1}}
\renewcommand{\maketitle}{%
  \twocolumn[{%
    \begin{center}
      {\Large\bfseries \@title\par}\vspace{0.4ex}
      {\normalsize \@author\par}\vspace{0.2ex}
      {\small \@affiliation\par}\vspace{0.2ex}
      {\small \@date\par}
    \end{center}\vspace{0.6ex}
  }]}
\makeatother

\title{Uncertainty-aware RUL prediction for predictive maintenance: Sketch of Idea}
\author{Tobias Dietze, Tobias Meier, Daria Samsonova, Srisharanya Vobblisetty}
\affiliation{DHBW Mannheim / Intelligence Engineering}
\date{30.11.2026}

\begin{document}
\maketitle

% =====================================================================
\section{Project Scope}
% Q: Is the project scope clearly defined?
To support a maintence decision an interval of remaining useful life (RUL) is easier to act on than a point estiamte.
This work will evaluate \textbf{(RQ1)}which uncertainty method gives the most reliable RUL intervals, and \textbf{(RQ2)} if calibrated intervals
do lead to better maintenance decisions than point predictions?

\ph{Concrete, measurable objective, e.g.\ ``classify X from sensor Y with $\geq$ Z\% macro-F1''.}

\ph{what we will and will not deliver.}

% =====================================================================
\section{Related Work}
Data driven RUL prediction is commonly benchmarked on the C-MAPSS dataset, which is made up of sensor readings from a simulated
turbofan~\cite{4711414}. Standard on this task are deep models which operate on sliding windows, such as CNNs~\cite{LI20181}.
However, these models produce point estimates. Basora et al. \cite{BASORA2025110513} benchmark different uncertainty quantification methods for deep prognostics,
including deep ensembles and MC dropout. Conformal prediction has been used to wrap RUL predictors in intervals with coverage guarantees \cite{javanmardi2023conformal,robinson2026riskaware}.
Mitici et al. \cite{MITICI2023109199} use probabilistic RUL prognostics for maintenance planning. To the knowledge of the authors, how the choice of uncertainty method affects the resulting maintenance
decisions has not been studied.

% =====================================================================
\section{Approach}
% Q: Is the choice of approach well founded?
\textbf{Method.} \ph{Model / algorithm and pipeline in brief (inputs $\rightarrow$ processing $\rightarrow$ outputs).}

\noindent\textbf{Rationale.} \ph{Why this approach fits the problem, data, and constraints (e.g.\ latency, memory, data size).}

\noindent\textbf{Alternatives considered.} \ph{Other options and why they were rejected.}

% =====================================================================
\section{Training \& Evaluation}
% Q: Are training and evaluation of the approach rigorous?
\begin{itemize}
  \item \textbf{Data:} \ph{source, size, labels, known biases / class imbalance.}
  \item \textbf{Splits:} \ph{train/val/test protocol; how leakage is prevented (e.g.\ split by subject/session/device).}
  \item \textbf{Training:} \ph{loss, optimiser, hyperparameter search, compute budget.}
  \item \textbf{Metrics:} \ph{primary and secondary metrics, and why they fit.}
  \item \textbf{Baselines:} \ph{trivial baseline plus at least one strong baseline.}
  \item \textbf{Rigor:} \ph{multiple seeds, confidence intervals / significance tests, ablations.}
\end{itemize}

% =====================================================================
\section{Reproducibility \& Deployment}
% Q: How are reproducibility and deployment handled?
\begin{itemize}
  \item \textbf{Code \& data:} \ph{repository URL, licence, data availability / access.}
  \item \textbf{Environment:} \ph{pinned dependencies (lock file), container, hardware used.}
  \item \textbf{Determinism:} \ph{fixed seeds, versioned configs, experiment tracking.}
  \item \textbf{Deployment:} \ph{target platform (e.g.\ edge device, server), runtime constraints, how the model is shipped and monitored.}
\end{itemize}

% =====================================================================
\section{AI Usage Statement}
% Q: Is an AI usage statement included?
\ph{State which AI tools were used (name and version), for what (e.g.\ code generation, writing, literature search), and how outputs were checked. If none were used, say so explicitly.}

% =====================================================================
{\small
  \bibliographystyle{abbrvnat}
  \bibliography{refs}
}

\end{document}
